\documentclass[11pt,a4paper]{article}
\usepackage[margin=2.2cm]{geometry}
\usepackage{amsmath,amssymb}
\usepackage{graphicx}
\usepackage{booktabs}
\usepackage[font=small,labelfont=bf]{caption}
\usepackage{xcolor}
\usepackage[colorlinks=true,linkcolor=blue!50!black,urlcolor=blue!50!black]{hyperref}
\usepackage{enumitem}
\setlist{itemsep=2pt,topsep=4pt}

\graphicspath{{../exp_data/}}
\newcommand{\fig}[1]{\includegraphics[width=\textwidth]{#1}}

\title{Why CG slows down on the cantilever}
\author{muFFTTO, \texttt{experiments/linear\_elastic\_mesh\_study}}
\date{23 October 2026}

\begin{document}
\maketitle

\noindent A linear elastic beam in a periodic cell with void pixels, solved by conjugate gradients (CG) with the Green
preconditioner. On finer meshes CG needs more and more iterations. These notes trace where that comes from and how to
remove it.

\medskip
\noindent\fbox{\parbox{\dimexpr\textwidth-2\fboxsep-2\fboxrule}{%
\textbf{Short answer.} Two separate things make CG depend on the mesh, and both can be removed.
\begin{itemize}
  \item \textbf{The void is too thin.} With a fixed number of void pixels, the void gets physically thinner every time
        the mesh is refined. That pushes the smallest eigenvalue towards zero. With a void of fixed \emph{physical}
        thickness, the smallest eigenvalue no longer depends on the mesh.
  \item \textbf{Dirichlet conditions (fixed nodes) spoil the largest eigenvalue.} A clamped face makes the largest
        eigenvalue grow like $1/h$. With forces only (Neumann conditions), the largest eigenvalue stays exactly 1.
\end{itemize}
With both fixes, the cantilever loaded only by forces needs 30--33 CG iterations on every mesh from $h = 1/40$ to
$1/640$. The clamped cantilever needs 40--94. The same holds in 3D (42 iterations up to 2 million voxels,
Section~\ref{sec:3d}). Fewer fixed nodes help only like a square root, and a stiff layer instead of the clamp works
only with a physical thickness and still costs more than the clamp (Sections~\ref{sec:nodes}, \ref{sec:stiff}).

\smallskip
\textbf{Dirichlet conditions can be kept.} The \emph{corrected Green} preconditioner is the exact inverse of the
reference problem with the fixed nodes (Section~\ref{sec:corrected}). With it, $\lambda_{\max} = 1$ for any fixed nodes,
and the clamped cantilever needs 30--32 CG iterations on every mesh, the same as the beam loaded by forces only.}}

%-----------------------------------------------------------------------------------------------------------------------
\section{The test problem}

A horizontal beam, length $L = 1$ and height $H = 0.2$, linear elastic ($E = 1$, $\nu = 0.3$, plane strain, small
strain), discretised with bilinear pixel elements. The beam sits in a periodic cell and is surrounded by void pixels
with zero stiffness, so it does not touch its periodic copies. The left face is clamped and a shear force acts on the
right face (Figure~\ref{fig:setup}).

The mesh is refined by changing the number of pixels over the beam height, $n_H$, from 2 to 256. The physical beam
stays the same.

\begin{figure}[h]
  \centering\includegraphics[width=0.75\textwidth]{cantilever_mesh_study_setup.png}
  \caption{The periodic cell at $n_H = 4$: void pixels (white), beam (yellow), clamped nodes (black squares) and the
  shear load on the right face (red).}
  \label{fig:setup}
\end{figure}

\paragraph{How CG is stopped.} CG stops when the estimated error in the energy norm, $\|u - u_k\|_K$, is below
$10^{-6}\,\|u\|_K$ (Hestenes--Stiefel lower estimate, minimal delay 10). This measures the error the way the elastic
energy sees it and does not depend on the preconditioner, so iteration counts can be compared between runs.

\paragraph{How the error is plotted.} The plots show the exact CG error, computed from the energy functional
$\Pi(v) = \tfrac12 v^\top K v - f^\top v$:
\[
  \|u - u_k\|_K^2 = 2\bigl(\Pi(u_k) - \Pi(u)\bigr), \qquad \|u\|_K^2 = -2\,\Pi(u).
\]
Since $r_k = f - K u_k$, $\Pi(u_k) = -\tfrac12 (f + r_k)^\top u_k$: two dot products per iteration.

%-----------------------------------------------------------------------------------------------------------------------
\section{What we saw: a long stall, then fast convergence}

On every mesh, the CG error stays at 100\,\% for many iterations, then drops quickly, at about one decade per 8--10
iterations (Figure~\ref{fig:cg}). The drop rate is the same on every mesh. Only the stall gets longer as the mesh gets
finer. With 2 void pixels, the number of CG iterations grows almost like $1/h$: 85, 136, 237 and 429 iterations for
$n_H = 32$, 64, 128 and 256.

\begin{figure}[h]
  \fig{cantilever_mesh_study_cg_convergence.png}
  \caption{Left: exact CG error (solid), the stopping estimate (dotted) and the discretisation error of each mesh
  (dashed). Middle: the residual, which first grows by about $100\times$ and is useless as a stopping test here.
  Right: estimate divided by exact error. Once CG converges, the estimate is within 1--6\,\% of the exact error.
  During the stall it badly underestimates, which would matter for a loose tolerance.}
  \label{fig:cg}
\end{figure}

%-----------------------------------------------------------------------------------------------------------------------
\section{The stall follows the beam length, not the fixed nodes}

When the mesh is refined, three things grow together: the number of fixed nodes, the beam length in pixels, and the
number of pixels overall. To separate them, the mesh was held fixed ($n_H = 64$) and one thing was changed at a time
(Figure~\ref{fig:plateau}).
\begin{itemize}
  \item[\textbf{A.}] Beam length $L = 0.5$, 1, 2, 4 with the same fixed nodes. The stall grows with the length, by
        about one iteration per 5 pixels of beam.
  \item[\textbf{B, C.}] $8\times$ more fixed nodes, or $4\times$ fewer. The stall changes by only 10--15\,\%.
  \item[\textbf{D.}] A beam of length 4 with the load moved close to the clamp. The stall stays the same as with the
        load at the tip, so the load--clamp distance is not what matters.
\end{itemize}

\begin{figure}[h]
  \fig{plateau_study_n_H64.png}
  \caption{Exact CG error at $n_H = 64$. Only panel A (the beam length) moves the stall substantially. In panel D all
  four load positions stall until about iteration 270.}
  \label{fig:plateau}
\end{figure}

%-----------------------------------------------------------------------------------------------------------------------
\section{The spectrum behind it}

The CG coefficients $\alpha_k$, $\beta_k$ give, almost for free, the eigenvalues that CG ``sees'': the eigenvalues
(Ritz values) of the Lanczos matrix
\[
  T_{kk} = \frac{1}{\alpha_k} + \frac{\beta_{k-1}}{\alpha_{k-1}}, \qquad
  T_{k,k+1} = \frac{\sqrt{\beta_k}}{\alpha_k}.
\]
On small meshes these were checked against a full dense eigenvalue computation of $M^{-1}K$
(Figure~\ref{fig:dense}). The smallest and largest values agree to all printed digits. Some small eigenvalues are never
seen by CG, because this load does not excite their mode shapes.

\begin{figure}[h]
  \centering\includegraphics[width=0.62\textwidth]{spectrum_check_dense.png}
  \caption{All eigenvalues from a dense computation (grey) and the Ritz values from the CG run (blue). The zero
  eigenvalues of nodes that touch only void, and of fixed nodes, are left out.}
  \label{fig:dense}
\end{figure}

\begin{figure}[h]
  \fig{spectrum_study.png}
  \caption{With 2 void pixels: the smallest eigenvalue falls like $h^3$ and the largest grows like $1/h$, so the
  condition number grows like $h^{-4}$. CG still does much better than $\sqrt{\kappa}$ would suggest, because the very
  small eigenvalues are few and isolated; the stall is CG finding them one at a time. Right: the number of eigenvalues
  below 0.05 doubles each time the beam length doubles (2, 4, 7, 14).}
  \label{fig:spectrum}
\end{figure}

\noindent So two different things are going on: something pulls the smallest eigenvalue down, and something pushes the
largest one up. The next two sections deal with each.

\begin{table}[h]
  \centering
  \begin{tabular}{rrrrrr}
    \toprule
    $n_H$ & $h$ & CG its & $\lambda_{\min}$ & $\lambda_{\max}$ & $\kappa$ \\
    \midrule
    16  & 1/80   & 59  & $4.0\times10^{-6}$ & 6.8 & $1.7\times10^{6}$ \\
    32  & 1/160  & 85  & $6.1\times10^{-7}$ & 13.4 & $2.2\times10^{7}$ \\
    64  & 1/320  & 136 & $8.4\times10^{-8}$ & 26.6 & $3.2\times10^{8}$ \\
    128 & 1/640  & 237 & $1.1\times10^{-8}$ & 53.2 & $4.8\times10^{9}$ \\
    256 & 1/1280 & 429 & $1.4\times10^{-9}$ & 106 & $7.5\times10^{10}$ \\
    \bottomrule
  \end{tabular}
  \caption{Clamped cantilever with 2 void pixels: CG iterations and the extreme eigenvalues of $M^{-1}K$.}
\end{table}

%-----------------------------------------------------------------------------------------------------------------------
\section{Fix 1: give the void a fixed physical thickness}

This follows from the earlier study of a strip in vacuum. The Green preconditioner $M$ is the inverse of the stiffness
of a cell that is solid \emph{everywhere}, void included. Each eigenvalue of $M^{-1}K$ is therefore a ratio of two
energies of the same displacement $v$:
\[
  \lambda(v) = \frac{\text{energy of } v \text{ in the real beam}}
                    {\text{energy of } v \text{ in a cell that is solid everywhere}}.
\]
A deformation pattern with wavelength $\ell$ along the beam fades into the void over a distance of about $\ell/(2\pi)$.
If the void is thinner than that, the pattern has to be squeezed into a thin gap, its reference energy becomes large,
and $\lambda$ becomes small. With 2 void pixels the void is $2h$ thick, so it gets thinner on every refinement and
$\lambda_{\min}$ falls.

The test: the void thickness $t$ was fixed in physical units ($0.025$ to $0.8$) and the mesh refined
(Figure~\ref{fig:void}). With a fixed $t$, the smallest eigenvalue is the same on every mesh. A thicker void gives a
larger eigenvalue and fewer iterations, levelling off at about $t = 0.4$--$0.8$. The solution does not depend on $t$,
since the void carries no load.

\begin{figure}[h]
  \fig{void_study.png}
  \caption{Fixed physical void thickness $t$ against 2 void pixels (dashed). Second panel: $\lambda_{\min}$ is flat in
  $h$ for every fixed $t$. Right: the tip deflection differs by $10^{-11}$ between void thicknesses.}
  \label{fig:void}
\end{figure}

\noindent The iterations still grow somewhat (55 $\to$ 71 per halving of $h$ at $t = 0.8$), and the largest eigenvalue
still doubles on each refinement. That is the second effect.

%-----------------------------------------------------------------------------------------------------------------------
\section{Fix 2: avoid Dirichlet conditions}

With the solid as the Green reference, every eigenvalue of $M^{-1}K$ should be at most 1, because the real material
is never stiffer than the reference. Fixed nodes break this. CG uses the preconditioner with the fixed nodes projected
out, $Q M Q$, and that is not the inverse of the reference problem with those nodes fixed, $(Q K_{\mathrm{ref}} Q)^{-1}$.
Next to each fixed node the mismatch acts at the scale of one pixel, which is why it grows like $1/h$.

The test: the same beam, a void of 0.4 on each side, and four ways of holding it (Figures~\ref{fig:bcsetup}
and~\ref{fig:bc}, Table~\ref{tab:bc}).

\begin{figure}[h]
  \centering\includegraphics[width=0.9\textwidth]{bc_setups.png}
  \caption{Supports and computed deformation at $n_H = 16$: clamped left face; pin + roller (3 fixed values); a free
  beam with a balanced load. For the free beam the rigid motion was removed before drawing. The fourth setting, the
  cantilever loaded by forces only, is described in Section~\ref{sec:traction}.}
  \label{fig:bcsetup}
\end{figure}

\begin{figure}[h]
  \fig{bc_study.png}
  \caption{Top left: CG iterations versus $h$. Top middle: the largest eigenvalue grows like $1/h$ only for the
  clamped face; it is exactly 1 for the free beam and for the cantilever loaded by forces. Top right: the smallest
  eigenvalue is constant in $h$ in every case, because the void is now thick enough. Bottom left: CG on the finest mesh
  ($n_H = 128$). Bottom middle: error of the compliance.}
  \label{fig:bc}
\end{figure}

\begin{table}[h]
  \centering
  \small\setlength{\tabcolsep}{4pt}
  \begin{tabular}{llcccc}
    \toprule
    Support & Fixed & CG its & $\lambda_{\max}$ & $\lambda_{\min}$ & Mesh- \\
            & values & ($h = 1/40 \to 1/640$) & & & independent \\
    \midrule
    Clamped left face    & $2(n_H + 1)$ & 40 $\to$ 94 & 5.3 $\to$ 79 ($\propto 1/h$) & $6.6\times10^{-4}$ & no \\
    Pin + roller         & 3            & 41 $\to$ 48 & 1.4 $\to$ 1.9 & $5.2 \to 4.0\times10^{-3}$ & nearly \\
    Free, balanced load  & 0            & 24 $\to$ 27 & 1.000 & $1.85\times10^{-2}$ & yes \\
    Cantilever by forces & 0            & 30 $\to$ 33 & 1.000 & $1.85\times10^{-2}$ & yes \\
    \bottomrule
  \end{tabular}
  \caption{Supports study, void thickness 0.4, $n_H = 8 \to 128$.}
  \label{tab:bc}
\end{table}

%-----------------------------------------------------------------------------------------------------------------------
\section{The cantilever without Dirichlet conditions}
\label{sec:traction}

The reactions of a cantilever are known from statics: a shear force $F$ and a bending moment $F L$. So the clamp can
be replaced by these reactions, applied as forces on the left face. With $\eta = y - H/2$ and $I = H^3/12$:
\begin{align*}
  \text{right face:}\quad & t_y(\eta) = \frac{3F}{2H}\Bigl(1 - \frac{4\eta^2}{H^2}\Bigr)
      && \text{(parabolic shear, total } F), \\
  \text{left face:}\quad  & t_y(\eta) = -\frac{3F}{2H}\Bigl(1 - \frac{4\eta^2}{H^2}\Bigr), \quad
                            t_x(\eta) = \frac{F L}{I}\,\eta
      && \text{(shear } {-F} \text{ and bending moment } F L).
\end{align*}
The tractions are integrated exactly onto the nodes (two-point Gauss per pixel), so the nodal forces balance to
round-off. Nothing is fixed. The solution is unique only up to a rigid motion, which CG handles without trouble because
the loads do not act on rigid motions.

This is the classical end-loaded beam of Timoshenko and Goodier, for which the 2D elasticity solution is known
exactly. Its compliance (twice the strain energy) is
\[
  F\cdot u = \frac{F^2 L^3}{3 E' I} + \frac{6 F^2 L}{5 \mu H}, \qquad E' = \frac{E}{1 - \nu^2} \ \text{(plane strain)},
\]
which is exactly the Timoshenko beam with shear coefficient $5/6$. The results:
\begin{itemize}
  \item \textbf{CG:} 30--33 iterations on every mesh, with $\lambda_{\max} = 1$ and $\lambda_{\min}$ constant.
  \item \textbf{Accuracy:} the compliance converges to the exact value as exactly $h^2$ (error $3.6\times10^{-5}$ at
        $n_H = 128$). The clamped version converges as $h^{1.76}$, because of the stress singularity at the corners of
        the clamp.
  \item \textbf{Difference from a real clamp:} the left face may warp slightly instead of staying flat. That changes the
        solution only near the left end (Saint-Venant's principle).
\end{itemize}
\textbf{Limit:} this only works when the reactions follow from statics alone (statically determinate supports:
cantilever, simply supported beam). For a beam clamped at both ends, the reactions are unknown in advance and Dirichlet
conditions are needed.

%-----------------------------------------------------------------------------------------------------------------------
\section{Check in 3D}
\label{sec:3d}

The same beam in 3D: length 1, cross-section $0.2 \times 0.2$, trilinear voxels, the same material, preconditioner and
stopping test. Two supports (clamped left face, or the cantilever loaded by forces only, with the shear spread
parabolically over the height and uniformly over the width), and two voids (2 voxels on each side, or a physical
thickness of 0.2). The meshes go from $n_H = 4$ to 32, the largest with $224 \times 96 \times 96 \approx 2$ million
voxels (Figure~\ref{fig:3d}, Table~\ref{tab:3d}).

\begin{figure}[h]
  \centering\includegraphics[width=0.85\textwidth]{cantilever_3D_study.png}
  \caption{3D beam: CG iterations, largest and smallest eigenvalue versus the voxel size, and CG on the finest mesh
  ($n_H = 32$). Red: forces only; blue: clamped face. Solid: void 0.2; dashed: 2 void voxels.}
  \label{fig:3d}
\end{figure}

\begin{table}[h]
  \centering\small\setlength{\tabcolsep}{4pt}
  \begin{tabular}{lccc}
    \toprule
    Case & CG its ($n_H = 4 \to 32$) & $\lambda_{\max}$ & $\lambda_{\min}$ \\
    \midrule
    Forces only, void 0.2      & 39 $\to$ 42  & 1.000 & $9.3\times10^{-3}$, constant \\
    Forces only, 2 void voxels & 38 $\to$ 54  & $\le 1$ & $1.1\times10^{-2} \to 5.4\times10^{-3}$, falls \\
    Clamped, void 0.2          & 49 $\to$ 94  & 1.9 $\to$ 12.7 ($\propto 1/h$) & $2.4\times10^{-4}$, constant \\
    Clamped, 2 void voxels     & 50 $\to$ 141 & 1.8 $\to$ 12.3 & $1.1\times10^{-4} \to 1.2\times10^{-6}$, falls \\
    \bottomrule
  \end{tabular}
  \caption{3D beam, $n_H = 4, 8, 16, 32$.}
  \label{tab:3d}
\end{table}

\noindent Both claims hold in 3D:
\begin{itemize}
  \item \textbf{The void:} with a physical void thickness, the smallest eigenvalue is the same on every mesh. With 2 void
        voxels it falls, fastest for the clamped beam (about $h^{2.5}$). In 3D even the beam loaded by forces only feels
        the thin void ($5.4\times10^{-3}$ instead of $9.4\times10^{-3}$ at $n_H = 32$).
  \item \textbf{The supports:} without Dirichlet conditions, $\lambda_{\max} = 1$ and CG needs 42 iterations from half a
        thousand to 2 million voxels. A clamped face makes $\lambda_{\max}$ double on every refinement, as in 2D.
\end{itemize}

%-----------------------------------------------------------------------------------------------------------------------
\section{Does the number of fixed nodes matter?}
\label{sec:nodes}

Yes, together with the pixel size. Here only a part of the left face is clamped: $k$ nodes, centred on the face, with
both components fixed. The void is 0.4 thick, so it plays no role. The same $k$ on two meshes separates ``how many fixed
nodes'' from ``how small the pixels are'' (Figure~\ref{fig:nodes}, left two panels, Table~\ref{tab:nodes}).

\begin{table}[h]
  \centering\small
  \begin{tabular}{rccc}
    \toprule
    $k$ & $n_H = 16$ & $n_H = 32$ & $n_H = 64$ \\
    \midrule
    3  & 3.9  & 4.6  & 5.3 \\
    5  & 5.4  & 6.7  & 7.9 \\
    9  & 7.7  & 9.9  & 12.1 \\
    17 & 10.2 & 14.6 & 18.7 \\
    33 & --   & 20.1 & 28.3 \\
    65 & --   & --   & 39.9 \\
    \bottomrule
  \end{tabular}
  \caption{Largest eigenvalue $\lambda_{\max}$ for $k$ fixed nodes on the left face.}
  \label{tab:nodes}
\end{table}

All points follow one law within about $\pm 5\,\%$:
\[
  \lambda_{\max} \approx 0.6\,\sqrt{k\, n_H}, \qquad\text{i.e.}\qquad \lambda_{\max} \propto \sqrt{k/h}.
\]
\begin{itemize}
  \item A fully clamped face has $k = n_H + 1$ nodes, which gives $\lambda_{\max} \propto 1/h$, as seen before.
  \item With a fixed number of nodes, $\lambda_{\max}$ still grows, like $h^{-1/2}$.
  \item Fewer fixed nodes help, but only like a square root. The CG count grows mildly with $k$ (46 to 71 at
        $n_H = 64$).
\end{itemize}

%-----------------------------------------------------------------------------------------------------------------------
\section{A stiff layer instead of the clamp}
\label{sec:stiff}

Idea: replace the clamped face by a layer of very stiff material left of the beam, $\kappa$ times stiffer than the beam,
and fix only the bottom and the top node of that layer. Those two nodes remove the rigid motion; the large reaction
forces that hold the layer in place come out of them. The void is again 0.4 thick.

\paragraph{One pixel thick: does not work.} A layer one pixel thick gets thinner on every refinement. Its bending
stiffness falls like $\kappa h^3$, so it clamps worse on finer meshes, and CG also gets slower
(Figure~\ref{fig:nodes}, right panels). At $n_H = 128$ the layer with $\kappa = 100$ is 29\,\% too soft and CG needs
236 iterations; with $\kappa = 10^4$ it is 11\,\% too soft and needs 968 iterations.

\begin{figure}[h]
  \fig{fixed_nodes_study.png}
  \caption{Left two panels: largest eigenvalue and CG iterations for $k$ fixed nodes on three meshes. Other panels: a
  stiff layer one pixel thick, held at two nodes: largest eigenvalue, CG iterations, how much softer than the real
  clamp it is, and CG at $n_H = 64$.}
  \label{fig:nodes}
\end{figure}

\paragraph{Fixed physical thickness: works as predicted, but costs more than the clamp.} A layer 0.05 thick (a quarter
of the beam height) keeps the same stiffness on every mesh (Figure~\ref{fig:stiff}, Table~\ref{tab:stiff}).
\begin{itemize}
  \item \textbf{It clamps equally well on every mesh.} With $\kappa = 100$ the beam is about 1\,\% too soft, on every
        mesh.
  \item \textbf{The eigenvalues form two clusters}, one around 0.1--1 from the beam and one up to about $\kappa$ from the
        layer. They barely move with $h$: $\lambda_{\max} \approx 1.5$--$2\kappa$, growing about 8\,\% per halving of $h$.
  \item \textbf{CG handles the stiff cluster.} The cost grows by about 40--60 iterations per decade of $\kappa$, far less
        than $\sqrt{\kappa}$ would suggest. The convergence curves show steps: CG resolves the stiff-layer eigenvalues a
        few at a time.
  \item \textbf{But it is more expensive than the clamp.} For a 1\,\% clamp error ($\kappa = 100$) CG needs 153
        iterations at $n_H = 128$, against 94 for the real clamp and 33 for the beam loaded by forces only.
\end{itemize}

\begin{figure}[h]
  \fig{stiff_layer_study.png}
  \caption{Stiff layer held at two nodes. Solid: layer 0.05 thick; dashed: one pixel thick. Top: largest and smallest
  eigenvalue and CG iterations versus $h$. Bottom: how much softer than the real clamp, CG at $n_H = 64$, and the Ritz
  values at $n_H = 64$ for each $\kappa$ (two clusters).}
  \label{fig:stiff}
\end{figure}

\begin{table}[h]
  \centering\footnotesize\setlength{\tabcolsep}{4pt}
  \begin{tabular}{rccc}
    \toprule
    $\kappa$ & CG its ($n_H = 16 \to 128$) & $\lambda_{\max}$ & too soft vs.\ clamp \\
    \midrule
    1      & 37 $\to$ 41   & 1.6 $\to$ 2.0      & 53 $\to$ 95\,\% (no real clamp) \\
    10     & 73 $\to$ 88   & 15 $\to$ 20        & 5.6 $\to$ 9.8\,\% \\
    100    & 115 $\to$ 153 & 154 $\to$ 196      & 0.6 $\to$ 1.0\,\% \\
    1000   & 158 $\to$ 226 & 1540 $\to$ 1960    & 0.06 $\to$ 0.10\,\% \\
    $10^4$ & 201 $\to$ 289 & 15400 $\to$ 19600  & 0.006 $\to$ 0.010\,\% \\
    \midrule
    clamped face & 46 $\to$ 94 & 10 $\to$ 79 & -- \\
    forces only  & 31 $\to$ 33 & 1.000       & -- \\
    \bottomrule
  \end{tabular}
  \caption{Stiff layer 0.05 thick, held at two nodes, void 0.4.}
  \label{tab:stiff}
\end{table}

\noindent\textbf{Why.} The Green preconditioner uses one reference material. With the beam as reference, the layer's
eigenvalues sit near $\kappa$; with the layer as reference, the beam's would sit near $1/\kappa$. One reference cannot
serve both materials. A preconditioner that adapts locally, such as Green-Jacobi (Green plus a diagonal scaling), should
pull the $\kappa$ cluster back towards 1. That is the next test.

%-----------------------------------------------------------------------------------------------------------------------
\section{Fix 3: the corrected Green for Dirichlet conditions}
\label{sec:corrected}

Section~6 showed that fixed nodes break the bound $\lambda_{\max} \le 1$. The reason is the way the preconditioner
treats them: CG uses Green with the fixed dofs cut out, and that is not the inverse of the reference problem with
those dofs fixed. The corrected Green is that inverse. It is built from the ordinary Green operator plus a small
correction on the fixed dofs, so it keeps Dirichlet conditions and still has $\lambda_{\max} \le 1$.

\paragraph{Notation.} $\mathbf{K}(\mathbb{C}) = \mathbf{B}^{\top}\mathbf{W}\mathbb{C}\mathbf{B}$ is the stiffness matrix of
the material $\mathbb{C}$, $\mathbf{K}_0 = \mathbf{K}(\mathbb{C}_0)$ that of the constant reference, and
$\mathbf{G}_0 = \mathbf{K}_0^{-1}$ on the mean-free vectors the Green operator (FFT, zero frequency omitted). The $m$
fixed (Dirichlet) dofs are selected by the projection $\mathbf{P} \in \mathbb{R}^{n\times n}$ (diagonal, 1 at the fixed
dofs and 0 elsewhere), $\mathbf{Q} = \mathbf{I} - \mathbf{P}$ projects onto the free dofs, and $\mathbf{N} \in \mathbb{R}^{n\times d}$ holds the $d$ rigid
translations, which $\mathbf{G}_0$ leaves out ($\mathbf{G}_0\mathbf{N} = \mathbf{0}$).

\subsection*{The idea without formulas}

\paragraph{What a preconditioner does.} In every iteration, CG has a residual $\mathbf{r}$: the forces that are not yet in
balance. The preconditioner answers the question \emph{``how would the structure move under these forces?''}, not
exactly, but for a simpler structure whose answer is cheap. The closer that simpler structure is to the real one, the
fewer iterations CG needs.

\paragraph{The simpler structure of Green.} Green answers the question for the \emph{reference solid}: the whole periodic
cell filled with the beam material, voids filled in, and \emph{nothing held}. For that structure the answer is cheap,
because every pixel is the same, so the FFT solves it exactly. As long as the real problem has no supports, the real
beam differs from the reference solid only by its missing material (the void). Missing material can only make the
structure softer, never stiffer. That is where $\lambda_{\max} \le 1$ comes from.

\paragraph{What goes wrong with fixed nodes (plain Green).} The real problem now has nodes that must not move. Plain
Green ignores them. It computes the motion of the free reference solid, in which the fixed nodes move like any other
node, and then simply sets the motion of the fixed nodes to zero afterwards. Picture a rubber sheet that you push with
your hand, look at its shape, and then pin a few points back to their original place by force. Next to every pinned
point the sheet now has a sharp dent, one pixel wide. A dent that is one pixel wide stores a lot of energy, and the
smaller the pixels, the sharper and more energetic the dent. To CG this looks like a far too stiff answer, and that is
exactly the $\lambda_{\max} \propto 1/h$ we measured. The more nodes are pinned this way, the more dents (the
$\sqrt{k}$ in Section~\ref{sec:nodes}).

\paragraph{What the corrected Green does instead.} It answers the right question: how does the reference solid move
when the fixed nodes are \emph{held while it deforms}? It gets this from Green in three steps, by superposition:
\begin{enumerate}
  \item \textbf{Let it move freely.} Apply Green to the forces $\mathbf{r}$. The fixed nodes move by some amount; call it
        their \emph{error}.
  \item \textbf{Find the holding forces.} Ask which forces $\boldsymbol{\lambda}$, applied at the fixed nodes only, move
        the fixed nodes back by exactly that error. For this one needs to know how a unit force at one fixed node
        moves the other fixed nodes. That is the small table $\mathbf{S}$ (one row and column per fixed dof). It does
        not depend on $\mathbf{r}$, so it is computed once, before CG starts, by applying Green to one unit force per
        fixed dof.
  \item \textbf{Move it again with the holding forces.} Apply Green once more, now to the forces $\mathbf{r}$ minus the
        holding forces. Because the system is linear, the fixed nodes now end up exactly at zero, and the rest of the
        solid has moved smoothly around them, without dents.
\end{enumerate}
The holding forces are what the supports of a real structure would provide: the \emph{reaction forces} of the
reference solid. The result is the exact motion of the reference solid with the same supports as the real problem.
The real beam then again differs from its reference only by missing material, so $\lambda_{\max} \le 1$ is back, for
any number and position of fixed nodes.

\paragraph{One more detail: the translation.} Green does not know where the structure sits: a rigid shift of the whole
cell costs no energy, and Green leaves it out (the zero frequency). With supports the shift is no longer free, so the
corrected Green has to determine it as well. That adds $d$ unknowns $\boldsymbol{\alpha}$ (one shift per direction)
to the small table, together with the condition that the forces in step 3 are in balance.

\subsection*{The same in formulas}

So far CG used
\[
  \mathbf{M}_{\mathrm{G}} = \mathbf{Q}\,\mathbf{G}_0\,\mathbf{Q} \;\neq\; \bigl(\mathbf{Q}\mathbf{K}_0\mathbf{Q}\bigr)^{-1}:
\]
$\mathbf{G}_0$ is the free motion of the reference solid, and the $\mathbf{Q}$ on the left sets the fixed dofs to zero
afterwards (the dents). The corrected Green is the exact inverse, $\mathbf{M}_{\mathrm{GD}} =
(\mathbf{Q}\mathbf{K}_0\mathbf{Q})^{-1}$ on the free dofs. Applying it to a residual $\mathbf{r}$ means solving the
reference problem with the fixed dofs,
\[
  \begin{bmatrix} \mathbf{K}_0 & \mathbf{P} \\ \mathbf{P} & \mathbf{0} \end{bmatrix}
  \begin{bmatrix} \mathbf{z} \\ \boldsymbol{\lambda} \end{bmatrix}
  =
  \begin{bmatrix} \mathbf{r} \\ \mathbf{0} \end{bmatrix},
\]
where $\boldsymbol{\lambda} = \mathbf{P}\boldsymbol{\lambda}$ are the reaction forces, nonzero only at the fixed dofs.
With the Green operator,
\[
  \mathbf{z} = \mathbf{G}_0\bigl(\mathbf{r} - \boldsymbol{\lambda}\bigr) + \mathbf{N}\boldsymbol{\alpha},
\]
and $\boldsymbol{\lambda}$ and the translation $\boldsymbol{\alpha}$ follow from $\mathbf{P}\mathbf{z} = \mathbf{0}$
(the fixed dofs do not move) and $\mathbf{N}^{\top}(\mathbf{r} - \boldsymbol{\lambda}) = \mathbf{0}$ (the
reference problem is solvable only for a load in equilibrium):
\[
  \begin{bmatrix} \mathbf{S} & \mathbf{P}\mathbf{N} \\ \mathbf{N}^{\top}\mathbf{P} & \mathbf{0} \end{bmatrix}
  \begin{bmatrix} \boldsymbol{\lambda} \\ -\boldsymbol{\alpha} \end{bmatrix}
  =
  \begin{bmatrix} \mathbf{P}\mathbf{G}_0\mathbf{r} \\ \mathbf{N}^{\top}\mathbf{r} \end{bmatrix},
  \qquad
  \mathbf{S} = \mathbf{P}\mathbf{G}_0\mathbf{P}.
\]
The three steps above are: $\mathbf{y} = \mathbf{G}_0\mathbf{r}$ (step 1; its error at the fixed dofs is
$\mathbf{P}\mathbf{y}$), the small solve for $\boldsymbol{\lambda}$ and $\boldsymbol{\alpha}$ (step 2), and
$\mathbf{z} = \mathbf{G}_0(\mathbf{r} - \boldsymbol{\lambda}) + \mathbf{N}\boldsymbol{\alpha}$ (step 3).
$\mathbf{S}$ is the Green operator seen only at the fixed dofs; on the range of $\mathbf{P}$ it is an $m\times m$
symmetric positive definite matrix, and the small system is solved only there ($m + d$ unknowns). Without
the translations this is the Woodbury formula $\mathbf{M}_{\mathrm{GD}} = \mathbf{G}_0 -
\mathbf{G}_0\mathbf{P}\mathbf{S}^{-1}\mathbf{P}\mathbf{G}_0$ (with $\mathbf{S}^{-1}$ the inverse on the range of
$\mathbf{P}$). The idea is the capacitance matrix method of fast
Poisson solvers on irregular domains (Buzbee, Dorr, George and Golub 1971; Proskurowski and Widlund 1976), used here as a
preconditioner.

\paragraph{Cost.} Setup, once: $m$ Green applications to assemble $\mathbf{S}$, and one factorization of the
$(m+d)\times(m+d)$ matrix. In every CG iteration: two Green applications instead of one ($\mathbf{G}_0\mathbf{r}$ and
$\mathbf{G}_0(\mathbf{r} - \boldsymbol{\lambda})$) and one small solve.

\paragraph{Why it works.} If the real material is nowhere stiffer than the reference, $\mathbb{C}_q \le \mathbb{C}_0$
(void included), then $\mathbf{Q}\mathbf{K}\mathbf{Q} \le \mathbf{Q}\mathbf{K}_0\mathbf{Q}$, and therefore
$\lambda_{\max}(\mathbf{M}_{\mathrm{GD}}\mathbf{K}) \le 1$, whatever the number and position of the fixed dofs and
whatever $h$. Only the preconditioner changes; the solution is the same.

\paragraph{Results.} The same beam, void 0.4, $n_H = 16 \to 128$ (Table~\ref{tab:corrected}, Figure~\ref{fig:corrected}).

\begin{table}[h]
  \centering\footnotesize\setlength{\tabcolsep}{4pt}
  \begin{tabular}{llccc}
    \toprule
    Supports & Preconditioner & CG its, $n_H = 16, 32, 64, 128$ & $\lambda_{\max}$ & $\lambda_{\min}$ \\
    \midrule
    clamped face          & plain Green     & 46, 55, 71, 94 & 10 $\to$ 79 ($\propto 1/h$) & $6.6\times10^{-4}$ \\
                          & corrected Green & 30, 31, 32, 32 & 1.000 & $2.70\times10^{-3}$ \\
    pin + roller          & plain Green     & 43, 45, 47, 48 & 1.54 $\to$ 1.93 & $4.8 \to 4.0\times10^{-3}$ \\
                          & corrected Green & 29, 30, 31, 31 & 1.000 & $1.85\times10^{-2}$ \\
    3 fixed nodes         & plain Green     & 41, 45, 46 ($n_H \le 64$) & 3.9 $\to$ 5.3 & \\
                          & corrected Green & 34, 36, 38 ($n_H \le 64$) & 1.000 & \\
    stiff layer $\kappa = 100$, 2 nodes & plain Green & 115, 133, 143, 153 & 154 $\to$ 196 & $6.2\times10^{-4}$ \\
                          & corrected Green & 96, 111, 123, 132 & 100.0 & $2.6\times10^{-3}$ \\
    stiff layer $\kappa = 100$, by forces & plain Green & 94, 111, 118, 123 & 100.0 & $1.7\times10^{-2}$ \\
    forces only           & plain Green     & 31, 32, 32, 33 & 1.000 & $1.85\times10^{-2}$ \\
    \bottomrule
  \end{tabular}
  \caption{Plain against corrected Green. The solutions are identical in all printed digits.}
  \label{tab:corrected}
\end{table}

\begin{figure}[h]
  \fig{dirichlet_fix_study.png}
  \caption{Plain Green (dashed) against corrected Green (solid). Top: CG iterations, largest and smallest eigenvalue
  against $h$. Bottom: CG convergence at $n_H = 128$; the largest eigenvalue for partial clamps of 3 and 17 nodes; CG
  iterations of the stiff layer against its stiffness $\kappa$.}
  \label{fig:corrected}
\end{figure}

\begin{itemize}
  \item \textbf{Every Dirichlet case becomes mesh-independent.} The clamped face needs 30--32 iterations on every mesh,
        with $\lambda_{\max} = 1$ and $\lambda_{\min}$ constant. That is the level of the beam loaded by forces only.
  \item \textbf{Fixed dofs that only stop the rigid motion cost nothing.} With the pin + roller, the spectrum is exactly
        that of the beam loaded by forces only ($\lambda_{\min} = 1.85\times10^{-2}$).
  \item \textbf{The stiff layer keeps $\lambda_{\max} = \kappa$.} The corrected Green removes the Dirichlet part exactly,
        but a material stiffer than the reference still gives eigenvalues up to $\kappa$. Holding the layer by forces
        changes little. With the corrected Green the stiff layer is not needed: the clamp itself is cheaper.
  \item \textbf{Limit.} The setup grows with the number of fixed dofs $m$; see Section~\ref{sec:large_m}.
\end{itemize}

\subsection*{A beam with holes, and the cost in time}
\label{sec:holes}

To check that the result does not depend on the simple solid beam, the same comparison was run on a beam with four
circular holes of radius 0.05 on its centre line (Figure~\ref{fig:holes}). The holes are fixed in physical units, like
the void.

\begin{figure}[h]
  \fig{holes_study.png}
  \caption{Beam with four holes, void 0.4. Top: the deformed structure; CG iterations; largest and smallest eigenvalue
  against $h$. Bottom: CG convergence at $n_H = 128$; wall time (setup + CG, and setup alone); time per CG iteration.}
  \label{fig:holes}
\end{figure}

\begin{table}[h]
  \centering\small
  \begin{tabular}{lcc}
    \toprule
    Supports and preconditioner & CG its, $n_H = 16, 32, 64, 128$ & $\lambda_{\max}$ \\
    \midrule
    clamped face, plain Green     & 60, 74, 92, 122 & 10 $\to$ 79 \\
    clamped face, corrected Green & 41, 42, 41, 42  & $\le 1$ \\
    pin + roller, plain Green     & 50, 51, 52, 54  & 1.5 $\to$ 1.9 \\
    pin + roller, corrected Green & 37, 37, 37, 36  & 1.000 \\
    forces only                   & 39, 40, 40, 40  & $\le 1$ \\
    \bottomrule
  \end{tabular}
  \caption{Beam with four holes: CG iterations. The solutions of plain and corrected Green are identical in all printed
  digits.}
  \label{tab:holes}
\end{table}

The holes make the beam softer, so every case needs about ten iterations more than the solid beam, but the picture is
the same: with the corrected Green every case is mesh-independent.

\paragraph{Time.} Clamped face, one thread, each run timed alone:

\begin{center}
\small
\begin{tabular}{lcccc}
  \toprule
  & \multicolumn{2}{c}{$n_H = 64$} & \multicolumn{2}{c}{$n_H = 128$} \\
  & plain & corrected & plain & corrected \\
  \midrule
  setup [s]                 & 0.30 & 0.71 & 1.3 & 5.7 \\
  time per CG iteration [ms] & 122 & 133 & 562 & 594 \\
  CG iterations             & 92 & 41 & 122 & 42 \\
  \textbf{total [s]}        & \textbf{11.5} & \textbf{6.2} & \textbf{69.9} & \textbf{30.6} \\
  \bottomrule
\end{tabular}
\end{center}

\begin{itemize}
  \item \textbf{One iteration costs only 6--9\,\% more,} not twice as much. A CG iteration is dominated by the product
        with the stiffness matrix (strains and stresses at all quadrature points); the second Green application, an FFT
        pair, is comparatively cheap.
  \item \textbf{The setup} ($m = 258$ Green applications and a dense inverse at $n_H = 128$) adds 4.4\,s.
  \item \textbf{Overall the corrected Green is 2.3 times faster} at $n_H = 128$, and the gap grows on finer meshes,
        because plain Green needs more iterations on every refinement while the corrected Green does not.
\end{itemize}

\subsection*{Heterogeneous material}
\label{sec:hetero}

The beam with holes was also given a smooth heterogeneous stiffness, fixed in physical units: the beam material is
multiplied by $\kappa^{-\varphi}$ with $\varphi = \tfrac12\bigl(1 + \sin(2\pi x/0.2 + 0.7)\sin(2\pi y/0.1 + 0.3)\bigr)$, so
the stiffness varies between $1/\kappa$ and 1 times the beam material, for contrasts $\kappa = 10$ and $100$
(Figures~\ref{fig:het10} and~\ref{fig:het100}). The Green reference stays the stiffest material, so the bound
$\lambda_{\max} \le 1$ of the corrected Green still applies.

\begin{figure}[h]
  \fig{holes_study_het10.png}
  \caption{Beam with holes and heterogeneous stiffness, contrast 10. Same panels as Figure~\ref{fig:holes}; the
  structure is coloured by the stiffness. All runs 6 in parallel, so the times are rough.}
  \label{fig:het10}
\end{figure}

\begin{figure}[h]
  \fig{holes_study_het100.png}
  \caption{As Figure~\ref{fig:het10}, contrast 100.}
  \label{fig:het100}
\end{figure}

\begin{table}[h]
  \centering\small
  \begin{tabular}{lcc}
    \toprule
    & \multicolumn{2}{c}{CG its, $n_H = 16, 32, 64, 128$} \\
    Supports and preconditioner & $\kappa = 10$ & $\kappa = 100$ \\
    \midrule
    clamped face, plain Green               & 104, 120, 149, 192 & 219, 254, 309, 384 \\
    clamped face, corrected Green           & 78, 81, 85, 85     & 166, 187, 200, 204 \\
    clamped face, corrected Green, $\mathbf{S}$ only & 89, 93, 94, 94 & 196, 217, 236, 241 \\
    pin + roller, plain Green               & 89, 94, 98, 100    & 203, 220, 239, 243 \\
    pin + roller, corrected Green           & 74, 76, 76, 76     & 158, 177, 198, 201 \\
    forces only                             & 75, 77, 82, 81     & 166, 181, 200, 200 \\
    \bottomrule
  \end{tabular}
  \caption{Beam with holes and heterogeneous stiffness: CG iterations.}
  \label{tab:hetero}
\end{table}

\begin{itemize}
  \item \textbf{The corrected Green behaves like the problem without Dirichlet conditions.} For both contrasts its
        iteration counts follow those of the beam loaded by forces only, and $\lambda_{\max} \le 1$. Plain Green with the
        clamp grows to 192 and 384 iterations at $n_H = 128$, about twice as many.
  \item \textbf{The contrast sets the level.} The soft material lowers $\lambda_{\min}$ roughly in proportion to $1/\kappa$
        ($7.2\times10^{-4}$ at $\kappa = 10$, $1.9\times10^{-4}$ at $\kappa = 100$ for the clamp with the corrected Green),
        so all cases need more iterations at $\kappa = 100$. $\lambda_{\min}$ itself is constant in $h$.
  \item \textbf{Mild growth on the coarse meshes at $\kappa = 100$} (166 $\to$ 200) also appears for the forces-only
        problem, and it stops between $n_H = 64$ and 128. It comes from the sampled material: on coarse meshes the
        pixels do not yet resolve the stiffness field. It is not caused by the supports.
  \item \textbf{$\mathbf{S}$ only costs 10--20\,\% more iterations here} (94 vs.\ 85, 241 vs.\ 204). Its $\lambda_{\min}$
        equals that of plain Green: the missing translation leaves one small eigenvalue that CG has to resolve.
  \item \textbf{Time} (rough, all runs in parallel): at $n_H = 128$ the corrected Green needs 93\,s against 266\,s for
        plain Green ($\kappa = 10$), and 213\,s against 369\,s ($\kappa = 100$).
\end{itemize}

\subsection*{The correction with $\mathbf{S}$ only}
\label{sec:s_only}

The translation part ($\mathbf{N}$, $\boldsymbol{\alpha}$ and the balance condition) makes the correction look more
complicated than it is. A simpler version keeps only $\mathbf{S}$:
\[
  \mathbf{S}\boldsymbol{\lambda} = \mathbf{P}\mathbf{G}_0\mathbf{r}, \qquad
  \mathbf{z} = \mathbf{G}_0\bigl(\mathbf{r} - \boldsymbol{\lambda}\bigr), \qquad
  \mathbf{S} = \mathbf{P}\mathbf{G}_0\mathbf{P}.
\]
The fixed dofs still end at zero, $\mathbf{P}\mathbf{z} = \mathbf{P}\mathbf{G}_0\mathbf{r} - \mathbf{S}\boldsymbol{\lambda} =
\mathbf{0}$. What is missing is the shift: $\mathbf{G}_0$ returns only vectors with zero mean, so every direction CG
receives has zero mean, while the true solution usually does not.

\paragraph{Why it still works here.} The mean is taken over \emph{all} nodes of the cell, including the nodes in the
void. Their displacements cost no energy, so they can take up any mean without changing the solution in the solid.
On the beam with holes, the version with $\mathbf{S}$ only gives the same solution in all printed digits and needs
44 iterations at $n_H = 16$ and 32, against 41--42 for the full correction (and 60--74 for plain Green). With
heterogeneous material it needs 10--20\,\% more iterations than the full correction (Table~\ref{tab:hetero}): the
missing translation leaves one small eigenvalue, at the level of plain Green's $\lambda_{\min}$.

\paragraph{When it would not work.} In a cell without void, for example a solid that fills the whole periodic cell and
is clamped on an internal plane, there are no free nodes to take up the mean, and the translation part is needed.
Around a structure in void, the version with $\mathbf{S}$ only is enough.

\subsection*{How thin may the void be? Large grids}
\label{sec:thin_void}

The void around the beam costs pixels: with $t = 0.4$ on each side the cell is $1.8 \times 1.0$ while the beam is
$1 \times 0.2$, so about 90\,\% of the pixels carry no material. The beam with holes was therefore solved with void
thicknesses from $t = 0.4$ down to $t = 0.00625$ (one pixel at $n_H = 32$), always fixed in physical units
(Figure~\ref{fig:thin_void}). The solution in the beam is the same for every $t$ (the compliances agree in all printed
digits), so only the cost changes.

\begin{figure}[h]
  \fig{void_holes_study.png}
  \caption{Beam with four holes, void thickness $t$ from 0.4 down to 0.00625 on each side. Left: CG iterations against
  $h$ (corrected Green solid, plain Green dashed, forces only dotted); the thinnest void is run up to $n_H = 1024$
  ($h = 1/5120$). Middle: smallest eigenvalue at $n_H = 128$. Right: wall time and number of pixels at $n_H = 128$
  (6 runs in parallel, so the times are rough).}
  \label{fig:thin_void}
\end{figure}

\begin{table}[h]
  \centering\footnotesize\setlength{\tabcolsep}{4pt}
  \begin{tabular}{rrcccccc}
    \toprule
    & & \multicolumn{2}{c}{CG its, $n_H = 16, 32, 64, 128$} & & \multicolumn{2}{c}{$\lambda_{\min}$ at $n_H = 128$} &
    total time [s] \\
    $t$ & pixels & plain Green & corrected Green & forces only & corrected & forces only & corr.\,/\,plain \\
    \midrule
    0.4     & 737\,280 & 60, 74, 92, 122  & 41, 42, 41, 42 & 39--40 & $2.5\times10^{-3}$ & $1.6\times10^{-2}$ & 31 / 70 \\
    0.2     & 344\,064 & 62, 75, 95, 125  & 41, 42, 41, 41 & 39--40 & $2.3\times10^{-3}$ & $1.7\times10^{-2}$ & 22 / 49 \\
    0.1     & 196\,608 & 64, 78, 99, 132  & 41, 42, 41, 41 & 39--40 & $1.9\times10^{-3}$ & $1.8\times10^{-2}$ & 16 / 28 \\
    0.05    & 135\,168 & 69, 84, 106, 141 & 42, 43, 42, 42 & 40--41 & $1.3\times10^{-3}$ & $1.5\times10^{-2}$ & 8 / 20 \\
    0.025   & 107\,520 & 73, 91, 115, 153 & 45, 45, 45, 45 & 43--44 & $8.0\times10^{-4}$ & $1.0\times10^{-2}$ & 6 / 17 \\
    0.0125  &  94\,464 & 82, 101, 124, 170 & 49, 50, 49, 49 & 48--49 & $4.5\times10^{-4}$ & $6.1\times10^{-3}$ & 6 / 17 \\
    0.00625 &  88\,128 & --, 115, 147, 193 & --, 58, 58, 58 & 57--58 & $2.4\times10^{-4}$ & $3.3\times10^{-3}$ & 6 / 16 \\
    \bottomrule
  \end{tabular}
  \caption{Beam with holes, clamped face: void thickness $t$ on each side. Pixels and times at $n_H = 128$; the time is
  setup + CG. $t = 0.00625$ is less than one pixel at $n_H = 16$ and was not run there.}
  \label{tab:thin_void}
\end{table}

\begin{itemize}
  \item \textbf{With the corrected Green every void thickness is mesh-independent.} Halving $t$ lowers
        $\lambda_{\min}$ roughly by a factor 1.8 below $t \approx 0.05$, and the iteration count rises slowly, from 42
        to 58 between $t = 0.4$ and $0.00625$. The forces-only beam behaves the same, iteration for iteration.
  \item \textbf{Plain Green is not}, for any $t$, because its $\lambda_{\max}$ grows like $1/h$ (Section~\ref{sec:nodes}).
  \item \textbf{A thin void is cheaper overall.} The thin void needs more iterations but far fewer pixels: at
        $n_H = 128$ the cell has 8 times fewer pixels with $t = 0.00625$ than with $t = 0.4$. With the corrected Green the
        fastest runs are $t = 0.0125$--$0.025$ (5--6\,s against 31\,s for $t = 0.4$).
\end{itemize}

\paragraph{Large grids.} The thinnest void, $t = 0.00625$, was run up to $n_H = 1024$: a grid of $5184 \times 1088$, 5.6
million pixels, 32 times finer than $n_H = 32$.

\begin{center}
\footnotesize\setlength{\tabcolsep}{4pt}
\begin{tabular}{rrcccccc}
  \toprule
  & & \multicolumn{3}{c}{CG iterations} & \multicolumn{2}{c}{$\lambda_{\min}$ / $\lambda_{\max}$} & \\
  $n_H$ & grid & plain & corrected & forces only & plain & corrected & fixed nodes \\
  \midrule
  32   & $162 \times 34$     & 115 & 58 & 58 & & $2.34\times10^{-4}$ / 0.999 & 33 \\
  128  & $648 \times 136$    & 193 & 58 & 57 & $2.9\times10^{-6}$ / 52  & $2.36\times10^{-4}$ / 0.999 & 129 \\
  256  & $1296 \times 272$   & 262 & 58 & 57 & $2.9\times10^{-6}$ / 103 & $2.37\times10^{-4}$ / 0.999 & 257 \\
  512  & $2592 \times 544$   & 370 & 58 & 57 & $2.9\times10^{-6}$ / 206 & $2.37\times10^{-4}$ / 0.999 & 513 \\
  1024 & $5184 \times 1088$  & \textbf{520} & \textbf{58} & 57 & $2.9\times10^{-6}$ / 412 & $2.37\times10^{-4}$ / 0.999 & 1025 \\
  \bottomrule
\end{tabular}

\vspace{6pt}
\begin{tabular}{rccc}
  \toprule
  $n_H$ & corrected Green: setup + CG [s] & plain Green [s] & speed-up \\
  \midrule
  256  & 9 + 27 = 36        & 115  & 3.2 \\
  512  & 61 + 96 = 157      & 606  & 3.9 \\
  1024 & 533 + 360 = \textbf{893} & \textbf{5462} & \textbf{6.1} \\
  \bottomrule
\end{tabular}
\end{center}

\begin{itemize}
  \item \textbf{The corrected Green needs 58 iterations on every mesh} from $n_H = 32$ to 1024. Its $\lambda_{\min}$ is
        constant to three digits and $\lambda_{\max} = 0.999$. Plain Green grows from 115 to 520 iterations, roughly like
        $h^{-1/2}$, as its $\lambda_{\max}$ doubles on each refinement ($\sqrt{\kappa} \propto h^{-1/2}$). Both give the
        same compliance, $5.198629\times10^{-6}$ at $n_H = 1024$.
  \item \textbf{The speed-up grows with the mesh}: 3.2, 3.9 and 6.1 times (4 runs in parallel, so the times are rough).
  \item \textbf{At $n_H = 1024$ the setup is 60\,\% of the time.} Assembling $\mathbf{S}$ takes $m = 2050$ Green
        applications of about 0.25\,s each. Its cost grows like $m$ times the grid size, about $4.5\times$ per refinement
        in 2D, faster than the CG part ($4\times$). One more refinement and the setup would dominate clearly. This is the
        2D version of the problem of the next paragraph, and the inner CG proposed there would remove most of it.
\end{itemize}

\subsection*{Many fixed dofs: a clamped wall in 3D}
\label{sec:large_m}

$\mathbf{S}$ has one row and one column per fixed dof. For a clamped face in 2D that is $m = 2(n_H+1)$, at most a few
hundred. A clamped wall in 3D has $(n+1)^2$ nodes with 3 dofs each, so $m$ grows like $h^{-2}$:

\begin{center}
\small
\begin{tabular}{rrrl}
  \toprule
  $n_H$ & fixed dofs $m$ & dense $\mathbf{S}$ & setup ($m$ Green applications) \\
  \midrule
  16 & 867    & 6 MB   & fine \\
  32 & 3\,267  & 85 MB  & about 1 hour on 2 million voxels \\
  64 & 12\,675 & 1.3 GB & not feasible \\
  \bottomrule
\end{tabular}
\end{center}

\noindent What breaks is the assembly of $\mathbf{S}$ ($m$ Green applications) and its dense factorization
($\sim m^3$), not the method. \emph{Applying} $\mathbf{S}$ stays cheap: $\mathbf{S}\mathbf{y} =
\mathbf{P}\mathbf{G}_0\mathbf{P}\mathbf{y}$ is one Green application. Ways to keep it general:
\begin{enumerate}
  \item \textbf{Do not assemble $\mathbf{S}$; solve $\mathbf{S}\boldsymbol{\lambda} = \mathbf{P}\mathbf{G}_0\mathbf{r}$ by an inner
        CG} in every outer iteration. One inner step is one Green application, about 6\,\% of an outer iteration, so
        15--20 inner steps cost about one extra outer iteration. Starting the inner CG from the $\boldsymbol{\lambda}$ of
        the previous outer iteration should cut this further. If the inner solve is not exact, the preconditioner
        changes slightly from one outer iteration to the next; that needs flexible CG or a tight inner tolerance.
  \item \textbf{Precondition the inner solve.} $\mathbf{S}$ is the flexibility of the reference solid seen at the wall;
        its condition number grows roughly like $1/h$. If the wall is a flat grid plane, $\mathbf{S}$ is nearly
        translation invariant along it, and an FFT in the plane of the wall would be a cheap preconditioner for it.
  \item \textbf{Use forces where statics allows it.} For statically determinate supports no $\mathbf{S}$ is needed at
        all (Section~\ref{sec:traction}).
\end{enumerate}
The inner-CG version is the natural next test, on the 3D clamped beam of Section~\ref{sec:3d} ($n_H = 8 \to 32$).

%-----------------------------------------------------------------------------------------------------------------------
\section{Practical rules}
\begin{enumerate}
  \item \textbf{Size the void in physical units, not in pixels.} Its thickness should be at least comparable to the
        wavelength of the softest deformation divided by $2\pi$. For this beam, 0.4 (twice the beam height) was enough.
        With the corrected Green a much thinner void also works: any fixed thickness is mesh-independent, and
        $t = 0.0125$--$0.025$ (a tenth of the beam height) was the fastest overall, because it saves many more pixels
        than it costs iterations (Section~\ref{sec:thin_void}).
  \item \textbf{With fixed nodes, use the corrected Green.} It keeps all eigenvalues at or below 1 for any Dirichlet
        conditions, and CG needs the same number of iterations on every mesh (Section~\ref{sec:corrected}). Without it,
        prefer forces to fixed nodes where statics allows it.
  \item \textbf{With plain Green, a few fixed values are cheap; fixed faces are not.} Pinning 3 values to remove rigid
        motion costs almost nothing. Clamping a face makes the largest eigenvalue grow like $1/h$.
  \item \textbf{Mesh-independent does not mean few iterations.} The smallest eigenvalue is set by the geometry. A more
        slender beam or a thinner void gives a smaller value, constant in $h$ but with more iterations.
  \item \textbf{Stop on the energy error, not the residual.} Here the residual first grows by $100\times$ before it
        falls.
  \item \textbf{The cost of fixed nodes grows like a square root.} With $k$ fixed nodes, $\lambda_{\max} \approx
        0.6\sqrt{k\,n_H}$.
  \item \textbf{A stiff layer needs a physical thickness.} A layer one pixel thick clamps worse and converges slower on
        every refinement. With a fixed thickness it behaves well, but with plain Green it is still more expensive than a
        clamped face.
\end{enumerate}

\subsection*{Open questions}
\begin{itemize}
  \item The corrected Green with many fixed dofs (a clamped wall in 3D): $\mathbf{S}$ applied without assembling it
        and solved by an inner CG, possibly preconditioned by an FFT in the plane of the wall
        (Section~\ref{sec:large_m}).
  \item How the smallest eigenvalue depends on the slenderness $L/H$. For the void thickness, Section~\ref{sec:thin_void}
        gives the trend for this beam ($\lambda_{\min}$ about halves when $t$ halves below $t \approx 0.05$), but not a
        rule for other geometries.
  \item Whether Green-Jacobi brings the stiff-layer cluster near $\kappa$ back towards 1, which would make the stiff
        layer cheaper than the clamp.
\end{itemize}

\subsection*{Files}
All in \texttt{experiments/linear\_elastic\_mesh\_study/}. Each mesh is solved once and stored in
\texttt{exp\_data/runs/}; the figures are rebuilt from the stored runs.
{\raggedright
\begin{description}[font=\ttfamily\small,style=nextline]
  \item[cantilever\_mesh\_study.py] the solver. Options: \texttt{N\_H}, \texttt{VOID} or \texttt{VOID\_THICKNESS},
        \texttt{BEAM\_LENGTH}, \texttt{CLAMP\_COLUMNS}, \texttt{CLAMP\_SPAN}, \texttt{LOAD\_X},
        \texttt{BC=clamp|simple|free|traction|stiff|stiffforce}, \texttt{STIFF\_FACTOR}, \texttt{STIFF\_THICKNESS},
        \texttt{PRECOND=green\_dirichlet} (the corrected Green), \texttt{DENSE\_SPECTRUM}.
  \item[run\_parallel.sh] mesh study, one process per mesh (Figures~\ref{fig:setup}, \ref{fig:cg}).
  \item[run\_plateau.sh, plateau\_study.py] what sets the stall (Figure~\ref{fig:plateau}).
  \item[spectrum\_study.py] Ritz values and the dense check (Figures~\ref{fig:dense}, \ref{fig:spectrum}).
  \item[run\_void.sh, void\_study.py] fixed physical void thickness (Figure~\ref{fig:void}).
  \item[run\_bc.sh, bc\_study.py, plot\_bc\_setups.py] supports study and the cantilever by forces
        (Figures~\ref{fig:bcsetup}, \ref{fig:bc}).
  \item[cantilever\_3D\_study.py, run\_3D.sh] the 3D check, runs in \texttt{exp\_data/runs3d/}
        (Figure~\ref{fig:3d}).
  \item[run\_fixed\_nodes.sh, fixed\_nodes\_study.py] number of fixed nodes and the one-pixel stiff layer
        (Figure~\ref{fig:nodes}).
  \item[run\_stiff\_layer.sh, stiff\_layer\_study.py] stiff layer of fixed thickness (Figure~\ref{fig:stiff}).
  \item[run\_dirichlet\_fix.sh, dirichlet\_fix\_study.py] the corrected Green and the stiff layer held by forces
        (Figure~\ref{fig:corrected}).
  \item[run\_hetero.sh] the beam with holes and heterogeneous stiffness (Figures~\ref{fig:het10}, \ref{fig:het100});
        \texttt{HETERO=}$\kappa$ in the solver, \texttt{holes\_study.py --hetero} $\kappa$ for the figure.
  \item[run\_holes.sh, holes\_study.py] the beam with holes, iterations and time (Figure~\ref{fig:holes});
        \texttt{PRECOND=green\_dirichlet\_S} is the version with $\mathbf{S}$ only.
  \item[run\_void\_holes.sh, run\_void\_holes\_thin.sh, run\_void\_large.sh, void\_holes\_study.py] void thickness on
        the beam with holes, and $t = 0.00625$ up to $n_H = 1024$ (Figure~\ref{fig:thin_void}).
  \item[report/corrected\_green.tex] a short stand-alone note on the corrected Green.
\end{description}
}

\end{document}
